\documentclass[../main.tex]{subfiles}
\graphicspath{{\subfix{../assets/}}}

\begin{document}
\section{Abstract}
This report details a use of applied mathematics in the calculation of various forces within a hypothetical vehicle.
We take randomised values of a multitude of characteristics and measurements to display the intuitive application and derivation of formulas that may be extended to more complex models that better reflect real-world scenarios in the industry.
We also undertake this report to showcase some potential considerations when making these calculations, particularly with respect to choices made in the inclusion and exclusion of certain parameters.
A delegation of responsibilities is included at the end for convenient attribution.
\newpage

\tableofcontents
\newpage

\section{Introduction}
Braking systems are no doubt a crucial element of the modern automotive, and their complex design is only made possible by the countless hours engineers and mathematicians dedicate to ensure the safety of everyone.
Specifically, we combine the experimental calculation of the coefficient of friction, $\mu$, of a brake pad against a steel bar with randomised vehicular characteristics to provide a comprehensive series of calculations to consider the stopping distance and braking force of the hypothetical vehicle at a predetermined speed.
We additionally calculate a parking brake system to illustrate these calculations within a more static setting.
Applied mathematics is a far more diverse field than as represented here, but still an invaluable took for engineers nonetheless.
\end{document}
